% Generated from locale/tr/content/intuitionistic-logic/semantics/topological-semantics.tex; do not hand-edit.


\olfileid[tr]{int}{sem}{top}

\olsection{Topolojik Semantik}

Sezgici mantığa bir semantik vermenin başka bir yolu, matematiksel topoloji
kavramını kullanmaktır.

\begin{defn}
  $X$ bir küme olsun. $X$ \emph{üzerinde bir topoloji}, aşağıdaki özellikleri
  sağlayan bir $\Top{O} \subseteq \Pow{X}$ kümesidir. $\Top{O}$'nun
  !!{element}s{}ına topolojinin \emph{açık kümeleri} denir. $X$ kümesi ile
  $\Top{O}$ birlikte bir \emph{topolojik uzay} diye adlandırılır.
  \begin{enumerate}
  \item Boş küme ile bütün uzay açıktır: $\emptyset$, $X \in \Top{O}$.
  \item Açık kümeler sonlu kesişimler altında kapalıdır: $U$, $V \in \Top{O}$
    ise $U \cap V \in \Top{O}$.
  \item Açık kümeler keyfî birleşimler altında kapalıdır: $U_i \in \Top{O}$
    ise ve bu tüm $i \in I$ için geçerliyse $\bigcup \Setabs{U_i}{i \in I}
    \in \Top{O}$.
  \end{enumerate}
\end{defn}

Açık kümeler topluluğu bağlamdan çıkarılabiliyorsa bir topoloji için $X$
yazabiliriz; yine de $X$'e ancak açık kümeler verildikten sonra topoloji
denebileceğine dikkat edin.

\begin{defn}
  Sezgici önerme mantığının bir \emph{topolojik modeli}, bir $\mModel{X} =
  \tuple{X, \Top{O}, V}$ üçlüsüdür; burada $\Top{O}$,~$X$ üzerinde bir
  topolojidir ve $V$, her önerme değişkenine $\Top{O}$ içinde bir açık küme
  atayan fonksiyondur.

  Bir topolojik model~$\mModel{X}$ verildiğinde $\Prop{X}{!A}$'yı tümevarımlı
  olarak şöyle tanımlayabiliriz:
  \begin{enumerate}
  \item $\Prop{X}{\lfalse} = \emptyset$
  \item $\Prop{X}{p} = V(p)$
  \item $\Prop{X}{!A \land !B} = \Prop{X}{!A} \cap \Prop{X}{!B}$
  \item $\Prop{X}{!A \lor !B} = \Prop{X}{!A} \cup \Prop{X}{!B}$
  \item $\Prop{X}{!A \lif !B} = \Interior{(X \setminus \Prop{X}{!A}) \cup
    \Prop{X}{!B}}$
  \end{enumerate}
  Burada $\Interior{V}$, bir $V \subseteq X$ kümesini onun \emph{iç
  bölgesine}, yani içerdiği bütün açık kümelerin birleşimine gönderen
  fonksiyondur. Başka bir deyişle,
  \[
  \Interior{V} = \bigcup \Setabs{U}{U \subseteq V \text{ ve } U \in \Top{O}}.
  \]
\end{defn}

Her kümenin iç bölgesi, açık kümelerin birleşimi olduğu için daima açıktır.
Dolayısıyla $\Prop{X}{!A}$ daima açık bir kümedir.

Topolojik semantik son derece soyut olsa da onu güdüleyebilecek bazı düşünme
yolları vardır. $X$'in !!{element}s{}ının ya da ``noktalarının'', ifadelerin
değerlendirilebildiği noktalar olduğunu varsayın. $!A$'nın doğru olduğu tüm
noktaların kümesi,~$!A$'nın dile getirdiği önermedir. Her nokta kümesi olası
bir önerme değildir; yalnızca $\Top{O}$'nun !!{element}s{}ı böyledir. $!A
\Entails !B$ ancak ve ancak $!B$,~$!A$'nın doğru olduğu her noktada doğruysa,
yani $\Prop{X}{!A} \subseteq \Prop{X}{!B}$ ise geçerlidir; bu tüm~$X$'ler
içindir.
Saçma ifade~$\lfalse$ hiçbir zaman doğru değildir; dolayısıyla
$\Prop{X}{\lfalse} = \emptyset$.

$!B \land !C$, $!B \lor !C$ ve $!B \lif !C$'nin dile getirdiği önermeler,
$!B$ ile~$!C$'nin dile getirdikleriyle nasıl ilişkili olmalıdır ki sezgici
bakımdan geçerli yasalar sağlansın, yani $!A \Proves !B$ ancak ve ancak
$\Prop{X}{!A} \subseteq \Prop{X}{!B}$ olsun? $\lfalse \Proves !A$ olmasını
her $!A$ için isteriz; bu sağlanır, çünkü $\emptyset \subseteq U$ olur ve bu
tüm $U$'lar için geçerlidir. $!B \land !C \Proves !B$ olduğundan $\Prop{X}{!B \land !C}
\subseteq \Prop{X}{!B}$ olmasını ve benzer biçimde $\Prop{X}{!B \land !C}
\subseteq \Prop{X}{!C}$ olmasını isteriz. $W \subseteq U$ ve $W \subseteq V$
koşullarını sağlayan en büyük küme $U \cap V$'dir. Tersine, $!B \Proves !B
\lor !C$ ve $!C \Proves !B \lor !C$; dolayısıyla $\Prop{X}{!B} \subseteq
\Prop{X}{!B \lor !C}$ ve $\Prop{X}{!C} \subseteq \Prop{X}{!B \lor !C}$
olmasını isteriz. En küçük~$W$ kümesi, $U \subseteq W$ ve $V \subseteq W$
koşullarını sağlayan $U \cup V$ kümesidir.

$\lif$ tanımı daha çetrefillidir: $!A \lif !B$,~$!A$ ile birleştirildiğinde
$!B$'yi mantıksal olarak gerektiren en zayıf önermeyi dile getirir. $!A \lif
!B$'nin $!A$ ile birleşince~$!B$'yi mantıksal olarak gerektirdiği, $(!A \lif
!B) \land !A \Proves !B $ ifadesinden açıktır. Dolayısıyla $\Prop{X}{!A \lif
!B}$, $\Prop{X}{!A \lif !B} \cap \Prop{X}{!A} \subseteq \Prop{X}{!B}$
koşulunu sağlayan en büyük açık küme olmalıdır; tanımımız buradan çıkar.

